\newpage
\section{Untermannigfaltigkeiten}
\label{4.8}

\begin{definition}{Untermanigfaltigkeit}
    Eine Teilmenge $M\subset\mathbb{R}^n$ heißt $d$-dimensionale \textbf{Untermannigfaltigkeit} der Klasse $C^l$ ($l\ge 1$, $d\le n-1$), wenn es zu jedem Punkt $a\in M$ eine offene Umgebung $U\subset\mathbb{R}^n$ und $l$-mal stetig differenzierbare Funktionen $g_1,\dots,g_{\,n-d}:U\to\mathbb{R}$ gibt, so dass gilt:

    \begin{enumerate}
        \item[(a)] $M\cap U=\{x\in U:\;g(x)=0\},$ wobei $g=(g_1,\dots,g_{\,n-d})$.
        \item[(b)] $\operatorname{Rang}Dg(a)=n-d.$
    \end{enumerate}

    Damit ist die Jacobimatrix $Dg$ gegeben durch
    $$
    Dg=
    \begin{pmatrix}
        \displaystyle\frac{\partial g_1}{\partial x_1} & \cdots & \displaystyle\frac{\partial g_1}{\partial x_n}\\[6pt]
        \vdots & & \vdots\\[4pt]
        \displaystyle\frac{\partial g_{\,n-d}}{\partial x_1} & \cdots & \displaystyle\frac{\partial g_{\,n-d}}{\partial x_n}
    \end{pmatrix},
    \qquad
    g=
    \begin{pmatrix}
        g_1\\[2pt]
        \vdots\\[2pt]
        g_{\,n-d}
    \end{pmatrix},
    \qquad g:U\to\mathbb{R}^{\,n-d}.
    $$

    D.\,h. $M$ lässt sich lokal als \textbf{Nullstellenmenge} von $(n-d)$ $C^l$-Funktionen von $n$ Variablen mit linear unabhängigen Gradienten schreiben.

    Die Bedingung $\operatorname{Rang}Dg(a)=n-d$ ist äquivalent dazu, dass die Gradienten $\nabla g_1,\dots,\nabla g_{\,n-d}$ im Punkt $a$ linear unabhängig sind.
\end{definition}

\begin{satz}{Untermanigfaltigkeit als Graph}
    Eine Teilmenge $M\subset\mathbb{R}^n$ ist eine $d$-dimensionale Untermannigfaltigkeit der Klasse $C^l$ genau dann, wenn zu jedem Punkt $a\in M$ nach evtl.\ Umnummerierung der Koordinaten

    sei
    $$
    a=\bigl(a_1,\dots,a_d,a_{d+1},\dots,a_n\bigr)^T=(a',a'')^T
    $$
    mit
    $$
    a':=(a_1,\dots,a_d)^T\in\mathbb{R}^d,\qquad
    a'':=(a_{d+1},\dots,a_n)^T\in\mathbb{R}^{\,n-d}.
    $$

    Es gibt offene Umgebungen $U'\subset\mathbb{R}^d$ von $a'$ und $U''\subset\mathbb{R}^{\,n-d}$ von $a''$ und eine $l$-mal stetig differenzierbare Abbildung
    $$
    \varphi:U'\to U'',\qquad x'\mapsto \varphi(x')=x'',
    $$
    so dass
    $$
    M\cap (U'\times U'')=\{(x',x'')\in U'\times U'':\ x''=\varphi(x')\}.
    $$

    Hier: $x'=(x_1,\dots,x_d)^T$ und $x''=(x_{d+1},\dots,x_n)^T$.
    \vspace{1em}

    D.\,h.\ $M$ läßt sich lokal als \emph{Graph} von $(n-d)$ $C^l$-Funktionen von $d$ Variablen schreiben.    
\end{satz}

\begin{satz}{Untermanigfaltigkeit als $d$-dimensionale Ebene}
    Sei $E_d\subset\mathbb{R}^n$ die $d$-dimensionale Ebene
    $$
        E_d := \{(x_1,\dots,x_n)^\top\in\mathbb{R}^n : x_{d+1}=\dots = x_n = 0\}.
    $$

    Eine Teilmenge $M\subset\mathbb{R}^n$ ist genau dann $d$-dimensionale Untermannigfaltigkeit der Klasse $C^l$, falls es zu jedem $a\in M$ offene Umgebungen
    \begin{enumerate}
        \item $U=U(a)\subset\mathbb{R}^n$, $V\subset\mathbb{R}^n$ und
        \item eine $C^l$-invertierbare Abbildung $F:U\to V$ gibt,
        
    \end{enumerate}
     so dass
    $$
        F(M\cap U) = E_d\cap V.
    $$
    D.\,h.\ durch eine lokale Koordinatentransformation der Klasse $C^l$ läßt sich $M$ zu einer \emph{$d$-dimensionalen Ebene} im $\mathbb{R}^n$ geradebiegen.

    \bigskip

    \footnotesize{* $C^l$-invertierbar: bedeutet $F$ ist $l$-mal stetig differenzierbar, $F$ bijektiv und $F^{-1}$ auch $l$-mal stetig differenzierbar.}    
\end{satz}

\begin{definition}{Reguläre Parametrisierung/ Immersion}
    Sei $W\subset\mathbb{R}^d$ offen.
    Eine stetig differenzierbare Abbildung $\Phi: W \to \mathbb{R}^n$ heißt \emph{Immersion} (oder \emph{reguläre Parametrisierung}), wenn
    $$
    \operatorname{Rang}D\Phi(w)=d,\qquad \forall w\in W,
    $$
    d.\,h. die Jacobi-Matrix
    $$
    D\Phi(w)=\Bigg(\frac{\partial\Phi_i}{\partial w_j}(w)\Bigg)_{i=1,\dots,n\;;\;j=1,\dots,d}
    =
    \begin{pmatrix}
    \frac{\partial\Phi_1}{\partial w_1} & \cdots & \frac{\partial\Phi_1}{\partial w_d} \\
    \vdots & \ddots & \vdots \\
    \frac{\partial\Phi_n}{\partial w_1} & \cdots & \frac{\partial\Phi_n}{\partial w_d}
    \end{pmatrix}
    \in\mathbb{R}^{\,n\times d},
    $$
    hat vollen (maximalen) Rang $d$, welcher Notwendigerweise $d \leq n$ ist.   
\end{definition}

\begin{satz}{Parametrisierungssatz}
    Eine Teilmenge $M\subset\mathbb{R}^n$ ist genau dann eine $d$-dimensionale Untermannigfaltigkeit der Klasse $C^\ell$, wenn es zu jedem Punkt $a\in M$
    \begin{itemize}
    \item eine offene Umgebung $U\subset\mathbb{R}^n$ von $a$,
    \item eine offene Teilmenge $W\subset\mathbb{R}^d$ und
    \item eine Immersion $\Phi:W\to\mathbb{R}^n$ der Klasse $C^\ell$
    \end{itemize}
    gibt, die $W$ homöomorph\footnote{Eine Funktion $f:M_1\to M_2$ heißt \emph{Homöomorphismus}, wenn $f$ bijektiv und stetig ist und $f^{-1}$ stetig ist.} auf $M\cap U$ abbildet.

    Der Homöomorphismus $\Phi:W\to (M\cap U)\subset\mathbb{R}^n$ heißt lokale \emph{Parameterdarstellung} oder \emph{Karte} der Untermannigfaltigkeit $M$.    
\end{satz}

\begin{definition}{Tangentialvektoren}
    Sei $M\subset\mathbb{R}^n$ eine Untermannigfaltigkeit und $a\in M$ ein Punkt.
    \vspace{1em}

    Ein Vektor $v\in\mathbb{R}^n$ heißt \emph{Tangentialvektor} an $M$ im Punkt $a$, wenn es eine stetig differenzierbare Kurve
    $$
    \gamma\colon\; ]-\varepsilon,\varepsilon[ \to M\subset\mathbb{R}^n,\qquad (\varepsilon>0),
    $$
    gibt mit
    $$
    \gamma(0)=a \quad\text{und}\quad \gamma'(0)=v.
    $$

    Der \emph{Tangentialraum} $T_aM$ an $M$ in $a$ besteht aus allen Tangentialvektoren an $M$ in $a$.  
\end{definition}

\begin{satz}{Tangentialraum}
    Sei $M\subset\mathbb{R}^n$ eine $d$-dimensionale Untermannigfaltigkeit und $a\in M$. Dann gilt:
    \begin{itemize}
        \item[a)] $T_aM$ ist ein $d$-dimensionaler Untervektorraum von $\mathbb{R}^n$.
        \item[b)] Sei $\Phi:W\to V\subset M$, ($W$ offen in $\mathbb{R}^d$, $V$ offen in $M$) eine Karte von $M$ und $c\in W$ ein Punkt mit $\Phi(c)=a$. Dann bilden die Vektoren
            $$
            \frac{\partial\Phi}{\partial w_1}(c),\dots,\frac{\partial\Phi}{\partial w_d}(c)
            $$
            eine \emph{Basis} von $T_aM$.
        \item[c)] Sei $U\subset\mathbb{R}^n$ eine offene Umgebung von $a$ und sei $g:U\to\mathbb{R}^{\,n-d}$ eine stetig differenzierbare Funktion mit
            $$
            M\cap U=\{x\in U:\; g(x)=0\}
            $$
            und
            $$
            \operatorname{Rang} Dg(a)=n-d.
            $$
            Dann gilt
            $$
            T_aM=\{v\in\mathbb{R}^n:\;\langle v,\nabla g_j(a)\rangle = 0\ \text{für } j=1,\dots,n-d\},
            $$
            wobei $g=(g_1,\dots,g_{n-d})$ und $\nabla g_j(a)$ den Gradienten von $g_j$ in $a$ bezeichnet.
    \end{itemize}   
\end{satz}

\begin{definition}{Normalenvektoren}
    Sei $M\subset\mathbb{R}^n$ eine $d$-dimensionale Untermannigfaltigkeit und $a\in M$.
    \vspace{1em}

    Ein Vektor $v\in\mathbb{R}^n$ heißt \emph{Normalenvektor} von $M$ in $a$, wenn er auf $T_aM$ senkrecht steht, d.\,h.
    $$
    (v,w)=0\qquad\text{für alle }w\in T_aM.
    $$

    Die Normalenvektoren von $M$ in $a$ bilden deshalb einen $(n-d)$-dimensionalen Untervektorraum $N_aM\subset\mathbb{R}^n$.

    Nach Satz 4.33 c) ist
    $$
    \nabla g_1(a),\dots,\nabla g_{\,n-d}(a)
    $$
    eine Basis von $N_aM$.    
\end{definition}